% Синтетический кейс: xelatex + fontspec + polyglossia (без fontenc/inputenc).
\documentclass[12pt]{article}
\usepackage{fontspec}
\setmainfont{DejaVu Serif}
\usepackage{polyglossia}
\setmainlanguage{russian}
\usepackage{amsmath,amsthm,unicode-math}
\setmathfont{DejaVu Math TeX Gyre}

\newtheorem{theorem}{Теорема}

\begin{document}

\section{Юникод}

\begin{theorem}\label{thm:x}
Теорема в документе для XeLaTeX.
\end{theorem}

Формула с юникодом в тексте — α и β:
\begin{equation}\label{eq:x}
  \alpha + \beta = \gamma .
\end{equation}
См.~\ref{thm:x} и~\eqref{eq:x}.

\end{document}
